\section{The Expected Number of Solution Clusters}
\label{sec:solution-components}
\label{sec:direct-identity}

\begin{mainidentityrestatement}
For every finite self-loop-free multigraph \(\Gin\),
\[
  \E_{\Phi}\,\solcomp{\rform{\Gin}}
  =
  \Ex{2^{\Kdef}\left(\frac34\right)^{\Tdef}} .
\]
\end{mainidentityrestatement}

The proof has three steps, one in each subsection. We first show that the
number of solution clusters equals an Euler characteristic
(Section~\ref{subsec:euler-characteristic}), then write its expectation as a
signed sum over sets \(\freeset\) of free variables
(Section~\ref{subsec:red-blue-expansion}), and finally show that the deletion
experiment computes this sum (Section~\ref{subsec:deferred-definition}).

\medskip
\noindent\textbf{Notation of this section.} Besides the symbols of
Section~\ref{subsec:notation}, this section uses two new symbols, the face
\(\face{\freeset}{\fixasg}\) of the Boolean cube and the edge factor
\(\edgefac{\freeset}{e}\).

\subsection{Euler characteristic}\label{subsec:euler-characteristic}

Let \(\fml\) be a 2-CNF formula whose variables are indexed by a finite set
\(\Vset\); from Section~\ref{subsec:red-blue-expansion} on, \(\Vset\) is the
vertex set of \(\Gin\). The solution set \(\Sol(\fml)\) is a set of vertices of
the Boolean cube \(\{0,1\}^{\Vset}\), and to define its Euler characteristic we
also record the higher-dimensional faces that lie inside it.

For \(\freeset\subseteq \Vset\) and an assignment
\(\fixasg:\Vset\setminus \freeset\to\{0,1\}\), let
\[
  \face{\freeset}{\fixasg}
  :=
  \{\text{all assignments that agree with \(\fixasg\) outside \(\freeset\)}\}.
\]
The variables in \(\freeset\) are free, so \(\face{\freeset}{\fixasg}\) has
\(2^{|\freeset|}\) assignments, and we call it a face of dimension
\(|\freeset|\). For example, a face of dimension \(0\) is a point, a face of
dimension \(1\) is a line segment, and a face of dimension \(2\) is a square.

Let \(\cplx{\fml}\) be the set of all faces \(\face{\freeset}{\fixasg}\) with
\[
  \face{\freeset}{\fixasg}\subseteq\Sol(\fml).
\]
In other words, a face is in \(\cplx{\fml}\) if and only if all its assignments
satisfy \(\fml\). Since every subface of such a face has the same property,
\(\cplx{\fml}\) is a finite cubical complex, and its \emph{Euler characteristic}
is
\begin{equation}\label{eq:euler-definition}
  \euler{\cplx{\fml}}
  :=
  \sum_{\face{\freeset}{\fixasg}\in\cplx{\fml}}(-1)^{|\freeset|}.
\end{equation}
In words, we count points, minus line segments, plus squares, minus cubes, and
so on.

\begin{lemma}[Euler characteristic of the solution graph]\label{lem:euler-characteristic}
For every 2-CNF formula \(\fml\),
\begin{equation}\label{eq:component-cancellation}
  \solcomp{\fml}
  =
  \euler{\cplx{\fml}}
  =
  \sum_{\freeset\subseteq \Vset}
  \ \sum_{\fixasg:\Vset\setminus \freeset\to\{0,1\}}
  (-1)^{|\freeset|}
  \1_{\{\face{\freeset}{\fixasg}\subseteq\Sol(\fml)\}}.
\end{equation}
\end{lemma}

The second equality is just the definition~\eqref{eq:euler-definition},
written as a sum over all faces. The first equality is proved by induction on
the number of variables, and the induction step uses the following slicing
formula.

Choose a variable \(\bvar{u}\) of \(\fml\), and let \(\fml_0\) and \(\fml_1\)
be the 2-CNF formulas on the other variables obtained by setting
\(\bvar{u}=0\) and \(\bvar{u}=1\) in \(\fml\) and simplifying. Thus
\(\Sol(\fml_0)\) is the set of assignments \(y\) of the other variables such
that \((y,\bvar{u}=0)\) satisfies \(\fml\), and similarly for
\(\Sol(\fml_1)\). An assignment \(y\) lies in
\(\Sol(\fml_0\wedge\fml_1)=\Sol(\fml_0)\cap\Sol(\fml_1)\) when it works for both
values of \(\bvar{u}\), and then it gives an edge of the solution graph of
\(\fml\) between \((y,0)\) and \((y,1)\). As always, we view these three sets
as graphs.

\begin{lemma}[Slicing the solution graph]\label{lem:slice-components}
For every 2-CNF formula \(\fml\) and every variable \(\bvar{u}\),
\begin{equation}\label{eq:slice-components}
  \solcomp{\fml}=\solcomp{\fml_0}+\solcomp{\fml_1}-\solcomp{\fml_0\wedge\fml_1}.
\end{equation}
\end{lemma}

\begin{proof}
Write \(I:=\Sol(\fml_0\wedge\fml_1)\). If we leave out all edges of the solution
graph that flip \(\bvar{u}\), what remains is two copies, \(\Sol(\fml_0)\) at
\(\bvar{u}=0\) and \(\Sol(\fml_1)\) at \(\bvar{u}=1\), with
\(\solcomp{\fml_0}+\solcomp{\fml_1}\) components in total. Now we add the edges
that flip \(\bvar{u}\), which are given by the points of \(I\). We will show
that each component of \(I\) joins its own pair of components, one of
\(\Sol(\fml_0)\) and one of \(\Sol(\fml_1)\), so that each component of \(I\)
lowers the number of components by exactly \(1\).

The key point is that two different components of \(I\) cannot lie in the same
component of \(\Sol(\fml_0)\). Suppose that \(y,y'\in I\), and that
\[
  y=z_0,z_1,\ldots,z_k=y'
\]
is a path in \(\Sol(\fml_0)\). For each \(z_i\), define a new assignment
\(z_i'\) by
\[
  z_i'(v):=
  \begin{cases}
    y(v),&y(v)=y'(v),\\
    z_i(v),&y(v)\ne y'(v),
  \end{cases}
\]
for every vertex \(v\ne u\). In words, on every variable where \(y\) and \(y'\)
agree we change \(z_i\) to their common value, and on every other variable we
keep the value of \(z_i\).

We first check that \((z_i',0)\) satisfies \(\fml\), and for this we take any
literal. If it is \(\bvar{u}\) or \(\neg\bvar{u}\), then \((z_i',0)\), \((y,0)\)
and \((y',0)\) give it the same value. If its variable is one where \(y\) and
\(y'\) agree, then \((z_i',0)\) gives it the same value as \((y,0)\) and
\((y',0)\). Otherwise \((z_i',0)\) gives it the same value as \((z_i,0)\), and
also the same value as one of \((y,0)\) and \((y',0)\). In every case, if the
literal is false under \((z_i',0)\), then it is false under at least two of
\[
  (z_i,0),\qquad (y,0),\qquad (y',0).
\]
Now suppose that a clause is false under \((z_i',0)\), so that all of its
literals are false there. The clause has at most two literals, and each of them
is false under two of the three assignments above, so one of the three
assignments makes all literals of the clause false. This is a contradiction,
because all three assignments satisfy \(\fml\).

We next check \((z_i',1)\) in the same way, using the three satisfying
assignments
\[
  (z_i,0),\qquad (y,1),\qquad (y',1).
\]
For a literal of a variable other than \(\bvar{u}\) the argument is the same,
and for the literal \(\bvar{u}\) or \(\neg\bvar{u}\), the last two assignments
give it the same value as \((z_i',1)\), so it is again false under two of the
three. Thus \((z_i',1)\) also satisfies \(\fml\), and \(z_i'\in I\).

The assignments \(z_i\) and \(z_{i+1}\) differ in one variable, and the change
above does not increase the number of variables in which they differ. Hence
\(z_i'\) and \(z_{i+1}'\) are equal or differ in one variable, and after
removing repeated points,
\[
  y=z_0',z_1',\ldots,z_k'=y'
\]
is a path in \(I\). Therefore one component of \(\Sol(\fml_0)\) meets at most
one component of \(I\), and by exchanging \(0\) and \(1\), the same holds for
\(\Sol(\fml_1)\).

As a result, different components of \(I\) join different pairs of
components, and each of them lowers the number of components by \(1\). This
proves~\eqref{eq:slice-components}.
\end{proof}

\begin{proof}[Proof of Lemma~\ref{lem:euler-characteristic}]
We use induction on the number \(|\Vset|\) of variables.

If \(|\Vset|=0\), there is only the empty assignment. If it satisfies \(\fml\),
then the solution graph has one vertex and \(\cplx{\fml}\) has one face, of
dimension \(0\), so \(\solcomp{\fml}=\euler{\cplx{\fml}}=1\). If it does not
satisfy \(\fml\), then the solution graph and \(\cplx{\fml}\) are both empty,
and both sides are \(0\).

Now let \(|\Vset|\ge1\), and assume the result for all 2-CNF formulas with
fewer variables. Choose a variable \(\bvar{u}\), let \(\fml_0\), \(\fml_1\) be
as above, and divide the faces of \(\cplx{\fml}\) into three types:
\begin{itemize}
\item faces with \(\bvar{u}=0\); they correspond one-to-one to the faces of
  \(\cplx{\fml_0}\), with the same dimension;
\item faces with \(\bvar{u}=1\); they correspond one-to-one to the faces of
  \(\cplx{\fml_1}\), with the same dimension;
\item faces in which \(\bvar{u}\) is free; they correspond one-to-one to the
  faces of \(\cplx{\fml_0\wedge\fml_1}\), with dimension larger by \(1\).
\end{itemize}
Since the dimension in the third type is larger by \(1\), the sign in the Euler
characteristic is reversed there, and therefore
\begin{equation}\label{eq:euler-slice}
  \euler{\cplx{\fml}}
  =
  \euler{\cplx{\fml_0}}
  +
  \euler{\cplx{\fml_1}}
  -
  \euler{\cplx{\fml_0\wedge\fml_1}}.
\end{equation}
By the induction hypothesis, the right-hand side equals
\(\solcomp{\fml_0}+\solcomp{\fml_1}-\solcomp{\fml_0\wedge\fml_1}\), and by
Lemma~\ref{lem:slice-components} this is \(\solcomp{\fml}\). This completes the
induction.
\end{proof}

\subsection{Red-blue expansion}\label{subsec:red-blue-expansion}

Now fix a self-loop-free graph \(\Gin=(\Vset,\Eset)\). Taking the expectation of
Lemma~\ref{lem:euler-characteristic} gives
\begin{equation}\label{eq:expected-faces}
  \E_{\Phi}\,\solcomp{\rform{\Gin}}
  =\sum_{\freeset\subseteq \Vset}\ \sum_{\fixasg:\Vset\setminus \freeset\to\{0,1\}}
  (-1)^{|\freeset|}
  \Prx[\Phi]{\face{\freeset}{\fixasg}\subseteq\Sol(\rform{\Gin})}.
\end{equation}

Fix a set \(\freeset\) in the outer sum, color the free vertices (those in
\(\freeset\)) blue and the fixed vertices (those in \(\Vset\setminus\freeset\))
red, and fix an assignment \(\fixasg\) of the red vertices. Whether the random
clause on an edge holds on the whole face \(\face{\freeset}{\fixasg}\) depends
on its signs and on \(\fixasg\), but the probability of this event depends only
on the colors of the two endpoints:
\[
\begin{array}{c|c|c}
\text{endpoints} & \text{meaning} & \text{probability}\\
\hline
\text{red--red} & \text{both variables are fixed} & 3/4\\
\text{red--blue} & \text{one variable is free} & 1/2\\
\text{blue--blue} & \text{both variables are free} & 0
\end{array}
\]
For an edge \(e\), let \(\edgefac{\freeset}{e}\) be the number in the last
column; thus \(\edgefac{\freeset}{e}\) is \(3/4\), \(1/2\) or \(0\) when \(e\)
has \(0\), \(1\) or \(2\) endpoints in \(\freeset\).

\begin{lemma}[Red-blue formula]\label{lem:red-blue-expansion}
For every finite self-loop-free multigraph \(\Gin=(\Vset,\Eset)\) with
\(\nvert\) vertices,
\begin{equation}\label{eq:red-blue}
  \E_{\Phi}\,\solcomp{\rform{\Gin}}=
  \sum_{\freeset\subseteq\Vset}
  2^{\nvert-|\freeset|}(-1)^{|\freeset|}
  \prod_{e\in \Eset}\edgefac{\freeset}{e}.
\end{equation}
\end{lemma}

\begin{proof}
Fix \(\freeset\). We first show that all \(2^{\nvert-|\freeset|}\) choices of
\(\fixasg\) in~\eqref{eq:expected-faces} give the same probability. Let
\(\fixasg'\) be another assignment of the red vertices, and in every clause
flip the sign of each red variable on which \(\fixasg\) and \(\fixasg'\)
differ. This map is a bijection on sign patterns, and it maps the event
\(\face{\freeset}{\fixasg}\subseteq\Sol(\rform{\Gin})\) onto the event
\(\face{\freeset}{\fixasg'}\subseteq\Sol(\rform{\Gin})\). Since the signs are
independent and uniform, the map does not change probabilities, and hence the
inner sum in~\eqref{eq:expected-faces} equals
\(2^{\nvert-|\freeset|}(-1)^{|\freeset|}\) times the probability for one fixed
\(\fixasg\).

For fixed \(\freeset\) and \(\fixasg\) the clauses on different edges are
independent, so this probability is \(\prod_{e\in\Eset}\edgefac{\freeset}{e}\),
and summing over \(\freeset\) gives~\eqref{eq:red-blue}.
\end{proof}

The Euler characteristic has now become a sum over all red-blue colorings, that
is, over all sets \(\freeset\). In the next step we show that each edge factor
\(\edgefac{\freeset}{e}\) is the average of a random edge deletion.

\subsection{Random edge deletion and deferred revelation}\label{subsec:deferred-definition}

We use the deletion experiment of Section~\ref{subsec:overview-main-identity},
with its forest \(\Ffor\) and its numbers \(\Tdef\) and \(\Kdef\).

\begin{lemma}[Random edge-deletion formula]\label{lem:deletion-identity}
For every finite self-loop-free multigraph \(\Gin=(\Vset,\Eset)\) with
\(\nvert\) vertices,
\[
  \Ex{2^\Kdef\left(\frac34\right)^\Tdef}
  =
  \sum_{\freeset\subseteq\Vset}
  2^{\nvert-|\freeset|}(-1)^{|\freeset|}
  \prod_{e\in \Eset}\edgefac{\freeset}{e}.
\]
\end{lemma}

Lemmas~\ref{lem:deletion-identity} and~\ref{lem:red-blue-expansion} together
give Theorem~\ref{thm:overview-main-identity}.

\begin{proof}
\emph{Step 1: one edge.} Fix a set \(\freeset\), flip an independent fair coin
for each edge, and let \(Y_e(\freeset)\) be the following random contribution
of the edge \(e\): if \(e\) is deleted, then \(Y_e(\freeset)=1\), and if \(e\)
is kept, then \(Y_e(\freeset)\) is \(1/2\), \(0\) or \(-1\) when \(e\) is
red--red, red--blue or blue--blue. Averaging over the coin gives
\[
\begin{array}{c|c}
\text{endpoints} & \E\,Y_e(\freeset)\\
\hline
\text{red--red}
&\displaystyle \frac12\left(1+\frac12\right)=\frac34\\[0.7em]
\text{red--blue}
&\displaystyle \frac12(1+0)=\frac12\\[0.7em]
\text{blue--blue}
&\displaystyle \frac12(1-1)=0
\end{array}
\]
These are the three numbers of the previous table, so
\(\E\,Y_e(\freeset)=\edgefac{\freeset}{e}\) for every edge \(e\), and since the
coins of different edges are independent,
\[
  \E\prod_{e\in \Eset}Y_e(\freeset)
  =\prod_{e\in \Eset}\edgefac{\freeset}{e}.
\]

\emph{Step 2: exchange the sum and the expectation.} We insert this into the
right-hand side of the lemma; since both the sum and the expectation are
finite,
\begin{equation}\label{eq:swap-sum-expectation}
  \sum_{\freeset} 2^{\nvert-|\freeset|}(-1)^{|\freeset|}
    \prod_{e\in \Eset}\edgefac{\freeset}{e}
  =\Ex{\sum_{\freeset}
      2^{\nvert-|\freeset|}(-1)^{|\freeset|}
      \prod_{e\in \Eset}Y_e(\freeset)}.
\end{equation}
Before the exchange we fix \(\freeset\) and average over the coins; after it we
fix the coins and sum over \(\freeset\). For fixed coins, a deleted edge has
factor \(1\) and can be ignored, while a kept red--blue edge has factor \(0\).
Hence a set \(\freeset\) contributes only if every component of the kept graph
is all red or all blue.

\emph{Step 3: one forest component.} Now we reveal the coins in the order of
the deletion experiment. After the first stage we know the forest \(\Ffor\),
but the coins of the deferred edges are not yet revealed, and every deferred
edge has both endpoints in the same component of \(\Ffor\). Fix the outcome of
the first stage, take one component \(C\) of \(\Ffor\), and let \(t(C)\) be the
number of deferred edges in \(C\).

The tree \(C\) has \(|V(C)|-1\) edges, and they are all kept. If \(C\) had
vertices of both colors, the tree path between two vertices of different colors
would contain a red--blue edge, and the term would be \(0\). So only two
colorings of \(C\) contribute: all red and all blue. Let \(W_C\) be the number
of deferred edges in \(C\) that are kept at the end; then
\(W_C\sim\Bin(t(C),1/2)\), and \(C\) contains \(|V(C)|-1+W_C\) kept edges.

If \(C\) is all red, each of its vertices gives a factor \(2\) and each kept
edge gives a factor \(1/2\), so the term is
\[
  2^{|V(C)|}\left(\frac12\right)^{|V(C)|-1+W_C}=2^{1-W_C}.
\]
If \(C\) is all blue, each vertex and each kept edge gives a factor \(-1\), so
the term is
\[
  (-1)^{|V(C)|}(-1)^{|V(C)|-1+W_C}=-(-1)^{W_C}.
\]
Now we average over the deferred coins of \(C\). If \(t(C)=0\), then
\(W_C=0\), and the sum of the two terms is \(2-1=1\). If \(t(C)\ge1\), then
\[
  \Ex{2^{1-W_C}-(-1)^{W_C}}
  =2\left(\frac{1+\frac12}{2}\right)^{t(C)}
    -\left(\frac{1-1}{2}\right)^{t(C)}
  =2\left(\frac34\right)^{t(C)}.
\]
This has a simple meaning edge by edge. In the all-red term, a deferred edge
gives \(1\) if deleted and \(1/2\) if kept, so \(3/4\) on average, while in the
all-blue term it gives \(1\) or \(-1\), so \(0\) on average. Thus one deferred
edge is enough to make the all-blue term vanish on average.

\emph{Step 4: all components.} Different components of \(\Ffor\) have different
deferred coins, and a contributing coloring chooses all red or all blue for
each component independently. The conditional expectation given the first
stage is therefore a product over the components. A component with
\(t(C)=0\) gives \(1\), there are exactly \(\Kdef\) components with
\(t(C)\ge1\), and \(\sum_C t(C)=\Tdef\), so the conditional expectation is
\[
  \prod_{C:\,t(C)\ge1}2\left(\frac34\right)^{t(C)}
  =2^\Kdef\left(\frac34\right)^\Tdef.
\]
Finally, the coins revealed in this order have the same joint law as coins
flipped all at the start, so taking the expectation over the first stage gives
the lemma.
\end{proof}
